

BSE\cite{salpeter1951relativistic} is a useful tool developed based on the many-body perturbation theory to describe the dynamic interactions of excited particles (\emph{e.g.} excitons) through the exchange of virtual particles such as photons, in solids and molecules. When applied in conjunction with the quaisparticle energies from $GW$ calculations, it gives excellent outcomes in reproducing the experimentally observed excitionic emission spectra. For instance, \cref{fig:GW_approximation_diagram}(e) compares the transverse dielectric constant for monolayer \ce{MoS2}, $\varepsilon_{2}^{\perp}(\omega)$.\cite{ramasubramaniam2012large} It can be seen that the hybrid HSE functional produces two emission peaks (just below 2 eV) that are slightly higher in photon energy compared to the experimental spectrum. $G_{0}W_{0}$-approximation, on the other hand, overestimates the HSE-gap, by $\sim 1$ eV. Further incorporating the electron-hole interactions  through BSE is found to be essential in correcting the HSE-gap and produces the excitionic peaks at photon enegies that match exactly with the experimental spectrum. Note that the two exciton peaks observed in  \ce{MoS2} are originated from the spin-orbit couplings induced by the heavy Mo atom, which is also important to be included in the simulation. 

The general form of BSE resembles that of a Dyson equation, which is a matrix eigenvalue equation to be solved for obtaining the two-particle wavefunction $\Psi$. Hence, it has the general form of $\Psi=G_0 G_0 K\Psi$, in which $G_0$ is the free-particle wavefunction and $K$ is the interaction kernel that takes into account both the Coulomb attraction and exchange repulsion between two particles. For modelling the excitonic interactions in solids, BSE takes the following form
\begin{equation}
    \label{eq:BSE}
    \left( E_c - E_v \right) A_{vc\bm{k}} + \sum_{v'c'\bm{k}'} K^{eh}_{vc\bm{k},v'c'\bm{k}'} A_{v'c'\bm{k}'} = \Omega A_{vc\bm{k}},
\end{equation}
which $E_{c}$ and $E_{v}$ refer to the quasiparticle conduction and valence band energies from $GW$ calculations, whereas the interaction kernel $K$ further takes into account the electronic states at different reciprocal space vectors through the $\bm{k}$ indicies. $A$ contains the information of the  amplitudes of electron excitation from the valence to the conduction band at wavevector $k$. Solving the BSE thus gives the excitation energy $\Omega$ that can be directly matched with the optical spectra. Solving the BSE is computationally intensive, which quadratically scales with the product of  the number of valence and conduction bands, as well as the $k$ points involved. As such, in the practical application of the BSE, one often restricts the computation to only involve  the conduction band, the valence band and a few other bands directly above/below.


